\documentclass[aspectratio=169, 10pt]{beamer}
\usepackage[utf8]{inputenc}
\usepackage{xcolor}
\usepackage{array}

% -----------------------------------------------------------------------------
% Nerd Mode Dark Palette
% -----------------------------------------------------------------------------
\definecolor{nerdDark}{HTML}{0D1117}      % Background canvas
\definecolor{nerdPanel}{HTML}{161B22}     % Card/Code block background
\definecolor{nerdGreen}{HTML}{00FF66}     % Terminal Matrix Green
\definecolor{nerdCyan}{HTML}{00F0FF}      % Electric Cyan
\definecolor{nerdRed}{HTML}{FF0088}       % Electric Red (Dominic)
\definecolor{nerdYellow}{HTML}{FFFF00}    % Plain Yello (Dominic)
\definecolor{nerdPurple}{HTML}{BD93F9}    % Neon Purple
\definecolor{nerdText}{HTML}{C9D1D9}      % High-contrast Off-White Text
\definecolor{nerdDim}{HTML}{484F58}       % Dimmed Gray

\setbeamercolor{background canvas}{bg=nerdDark}
\setbeamercolor{normal text}{fg=nerdText}
\setbeamercolor{title}{fg=nerdCyan}
\setbeamercolor{author}{fg=nerdText}
\setbeamercolor{date}{fg=nerdDim}
\setbeamertemplate{navigation symbols}{}

% -----------------------------------------------------------------------------
% Metadata
% -----------------------------------------------------------------------------
\title{\textbf{SimHash: Schnelle Detektion von Plagiaten}}
\author{Dominic van der Zypen}
\date{}

\begin{document}

% =============================================================================
% SLIDE 1: Title Slide
% =============================================================================
\begin{frame}
  \titlepage
  \vspace*{-3cm}
  \begin{center}{2026-09-07}\end{center}
\end{frame}

% =============================================================================
% SLIDE 2: Was ist ein Hash?
% =============================================================================
\begin{frame}[fragile]
  \vspace{-0.2cm}
  % Title Header
  {\large \textcolor{nerdCyan}{\textbf{Was ist eine Hash-Funktion?}}}
  
  \vspace{0.15cm}

  \begin{itemize}
    \item Fixe Output-Bitl\"ange \textcolor{nerdGreen}{$n\in\mathbb{N}$}. 
    \item Ein \textcolor{nerdGreen}{{\bf Hash}} ist eine 
      Funktion $$h:\{0,1\}^*\to \{0,1\}^n.$$
    \item Beispiele: 
      \begin{enumerate}
        \item Initialen bei Namen.
        \item "Parity bit" + Verallgemeinerungen.
      \end{enumerate}
    \item Gew\"unschte Eigenschaft: \textcolor{nerdRed}{{\bf Avalanche effect}},
      d.h. kleine \"Anderungen im Input geben grosse \"Anderungen
      im Output.
  \end{itemize}
  
\end{frame}

% =============================================================================
% SLIDE 3: SimHash Demo Slide
% =============================================================================
\begin{frame}[fragile]
  \vspace{-0.2cm}
  % Title Header
  {\large \textcolor{nerdCyan}{\textbf{Beispiel f\"ur einen SimHash-Fingerprint}}}
  
  \vspace{0.15cm}
  
  \small
  \textcolor{nerdText}{\textbf{Input:}} \texttt{"Hallo Welt!"} \quad$\vert$\quad 
  \textcolor{nerdText}{\textbf{Anzahl Zeichen:}} \texttt{11} \quad$\vert$\quad 
  \textcolor{nerdText}{\textbf{Anzahl Trigramme:}} \texttt{9} \quad$\vert$\quad 
  \textcolor{nerdText}{\textbf{Mehrheit:}} $\ge \textcolor{nerdCyan}{5}$

  \vspace{0.25cm}
  
  % Table Container Panel
  \begin{center}
  \colorbox{nerdPanel}{%
  \ttfamily\small
  \begin{tabular}{l l l}
    \textcolor{nerdCyan}{\textbf{Trigramm}} & \textcolor{nerdCyan}{\textbf{32-bit-Hash (bin\"ar)}} & \textcolor{nerdCyan}{\textbf{Hexadezimal}} \\
    \hline
    "Hal" & \textcolor{nerdText}{1010 1001 0111 0101 0000 1111 1001 0110} & \textcolor{nerdPurple}{0xa9750f96} \\
    "all" & \textcolor{nerdText}{1010 1001 0000 1010 1100 1111 1000 0011} & \textcolor{nerdPurple}{0xa90acf83} \\
    "llo" & \textcolor{nerdText}{0011 0111 1110 0010 1101 1000 0100 0001} & \textcolor{nerdPurple}{0x37e2d841} \\
    "lo " & \textcolor{nerdText}{1100 1111 1111 0001 0011 1101 1110 0111} & \textcolor{nerdPurple}{0xcff13de7} \\
    "o W" & \textcolor{nerdText}{1100 1101 0011 1011 0111 0100 0101 1111} & \textcolor{nerdPurple}{0xcd3b745f} \\
    " We" & \textcolor{nerdText}{1011 0010 0001 0011 0010 1010 1010 0011} & \textcolor{nerdPurple}{0xb2132aa3} \\
    "Wel" & \textcolor{nerdText}{0000 1101 0001 1111 1101 0000 0111 1010} & \textcolor{nerdPurple}{0x0d1fd07a} \\
    "elt" & \textcolor{nerdText}{1111 1001 1010 1111 1010 0011 0001 1011} & \textcolor{nerdPurple}{0xf9afa31b} \\
    "lt!" & \textcolor{nerdText}{1001 0101 0100 0101 0110 1110 1110 1111} & \textcolor{nerdPurple}{0x95456eef} \\
    \hline
    \textcolor{nerdCyan}{\textbf{1-Count}} & \textcolor{nerdText}{[7, 8, 4, 4, 4, 4, 5, 5, 4, 5, 5, 6, 4, 5, 5, 4, \dots]} & \\
    \hline\hline
    \textcolor{nerdGreen}{\textbf{Fingerprint}} & \textcolor{nerdGreen}{\textbf{1010 1101 0011 0011 0110 1110 1100 0011}} & \textcolor{nerdGreen}{\textbf{0xad336ec3}} \\
  \end{tabular}%
  }
  \end{center}

  \vspace{0.1cm}
  \footnotesize
  \textcolor{nerdText}{\textbf{Regel:}} F\"ur jede Bit-Position $j$ gibt es eine "Abstimmung": Der Fingerprint ist 
  \textcolor{nerdGreen}{\texttt{1}}, falls $\text{voting\_vector}[j] \ge 5$, \\
  und \textcolor{nerdText}{\texttt{0}} sonst.
\end{frame}

% =============================================================================
% SLIDE 4: \"Ahnlichkeitsvergleich via Hamming-Distanz
% =============================================================================
\begin{frame}[fragile]
  \vspace{-0.2cm}
  % Title Header
  {\large \textcolor{nerdCyan}{\textbf{\"Ahnlichkeitsvergleich via Hamming-Distanz}}}
  
  \vspace{0.15cm}

  \begin{itemize}
    \item Die \textcolor{nerdGreen}{\textbf{Hamming-Distanz}} $d_H(x, y)$ 
	    z\"ahlt die Anzahl abweichender Bits:
      $$d_H(x, y) = \text{popcount}(x \oplus y)$$
    \item Kein Avalanche effect bei SimHash:\\
      \textcolor{nerdGreen}{\textbf{\"Ahnliche Dokumente $\Rightarrow$ kleine Hamming-Distanz.}}
  \end{itemize}

  \vspace{0.2cm}

  % Comparison Table Container
  \begin{center}
  \colorbox{nerdPanel}{%
  \ttfamily\small
  \begin{tabular}{l l l}
    \textcolor{nerdCyan}{\textbf{Dokument}} & \textcolor{nerdCyan}{\textbf{SimHash Fingerprint (32-Bit)}} & \textcolor{nerdCyan}{\textbf{Hex}} \\
    \hline
    \textcolor{nerdText}{Text A ("Hallo Welt!")} & 
	  \texttt{1010 1101 0011 0011 0110 1110 
	  1100 0011} & \textcolor{nerdPurple}{0xad336ec3} \\
    \textcolor{nerdText}{Text B ("Hallo Welt.")} & 
	  \texttt{1010 1101 0011 0011 0110 1110 
	  110\textcolor{nerdRed}{1} 00\textcolor{nerdRed}{0}1} & 
	  \textcolor{nerdPurple}{0xad336ed1} \\
    \hline
    \textcolor{nerdGreen}{\textbf{Bit-Vergleich}} & 
	  \texttt{\textcolor{nerdGreen}{0000 0000 0000 0000 0000 0000 
	  000\textcolor{nerdRed}{1} 
	  00\textcolor{nerdRed}{1}0}} & 
	  \textcolor{nerdYellow}{\textbf{2 Bits Diff}} \\
  \end{tabular}%
  }
  \end{center}

  \vspace{0.15cm}
  \footnotesize
  \textcolor{nerdText}{\textbf{Fazit:}} Da $d_H(A, B) = \textcolor{nerdGreen}{\mathbf{2}} \le k$ (typischer Schwellenwert $k \approx 3$), werden die Dokumente als \textcolor{nerdGreen}{\textbf{\"ahnlich / Plagiat}} eingestuft.
\end{frame}

\end{document}
